% Chapter 2 --- Contextual Inquiry

\chapter{Contextual Inquiry}\doublespacing
\label{Chapter2}
\lhead{Chapter II. \emph{Contextual Inquiry}}

%----------------------------------------------------------------------------------------
\section{Field Methodology and Site Observations}
%----------------------------------------------------------------------------------------

The fieldwork was conducted over several visits to procurement sites: a wheat procurement warehouse in active operation and a downstream processing factory. The method followed the contextual inquiry tradition \cite{Holtzblatt1998}. Researchers were embedded in the workspace, observing ongoing work and asking questions as they arose. This approach reveals how a process actually works, rather than how it is supposed to work.

Both sites presented demanding conditions. Warehouse assessment operates under external pressure: truck queues, noise, limited lighting control, and fine dust on every surface. The factory environment is marginally more controlled. Both share the challenge of sustaining precise visual judgement across a full working shift.

%----------------------------------------------------------------------------------------
\section{Stakeholders Involved}
%----------------------------------------------------------------------------------------

Four stakeholder groups were identified, each with a distinct relationship to the assessment process.

\textbf{Warehouse and factory assessors} perform the task directly. Their primary concern is throughput, because truck queues make slowness immediately visible. Under demanding shift conditions, assessors develop coping strategies for borderline specimens rather than deliberating on every ambiguous grain.

\textbf{Quality managers and administrators} interpret results and intervene when values approach rejection thresholds. They hold institutional memory about which vendors routinely supply borderline grain. Their judgements incorporate context that the formal protocol does not capture.

\textbf{Vendors and truck drivers} receive only the outcome: accept or reject. A vendor disputing a rejection has no access to the evidence or reasoning behind it.

\textbf{Food analysts} at headquarters define the quality standard itself. They set thresholds and maintain the category taxonomy. As Polanyi's analysis of expertise predicts \cite{Polanyi1966}, much of their knowledge is not codified in the training materials that warehouse assessors receive.

%----------------------------------------------------------------------------------------
\section{The ``As-Is'' Workflow: Periphery and Master Sampling}
%----------------------------------------------------------------------------------------

The current workflow uses a two-stage logic: rapid gatekeeping first, deeper confirmation second.

\subsection{Periphery Sample Testing}

When a truck arrives, grain portions are collected from outer \textit{boris} (gunny bags) to form a periphery sample. Three checks follow: moisture reading through a digital meter, Hectolitre measurement for variety validation (Dara, Lokwan, Sharbati, Amravati, Durum), and manual classification of a 100-gram sub-sample into twelve categories.

\begin{figure}[H]
    \centering
    \begin{subfigure}[b]{0.48\textwidth}
        \centering
        \includegraphics[width=0.8\textwidth]{moisture measurement in godown.jpg}
        \caption{Moisture content reading.}
        \label{fig:moisture_meter}
    \end{subfigure}
    \hfill
    \begin{subfigure}[b]{0.48\textwidth}
        \centering
        \includegraphics[width=0.8\textwidth]{Hectolitre weight measurement in godown.jpg}
        \caption{Hectolitre weight recording.}
        \label{fig:hectolitre_weight}
    \end{subfigure}
    \caption{Initial checks grains and impurities segregation: (a) moisture measurement and (b) Hectolitre weight for wheat variety verification.}
    \label{fig:periphery_tests}
\end{figure}

\subsection{Threshold Logic and the $W/W\%$ Calculation}

For each category, $W/W\%$ is calculated and compared with standard procurement thresholds. Acceptance requires impurity and defect percentages to remain below category limits. Pure wheat has no upper cap; a higher proportion indicates better lot quality.

\subsection{Master Sample and Redundancy}

If the periphery result is acceptable, unloading proceeds and a master sample is drawn from inner \textit{boris}. The full test is repeated to verify that inner-lot quality matches the outer-lot sample and to reduce adulteration risk.

\begin{figure}[H]
    \centering
    \includegraphics[width=0.45\textwidth]{factory sample collection from bori.jpeg}
    \caption{Master sample collection from \textit{boris} at godown.}
    \label{fig:sample_collection}
\end{figure}

%----------------------------------------------------------------------------------------
\section{Synthesis of Design Constraints}
%----------------------------------------------------------------------------------------

Four design constraints emerged from fieldwork:

\textbf{Latency:} Sixty minutes per sample is a physical consequence of one person sorting 2,800 grain entities by hand. The target cycle time for Drishti was set at under five minutes. Anything slower than ten minutes would provide insufficient operational benefit to justify adoption.

\textbf{Expertise concentration:} Experienced and inexperienced assessors are not interchangeable in practice. Agreement was high on major categories but varied significantly on borderline specimens, such as a kernel that reads as either \textit{Broken} or \textit{Sharbati} depending on assessor experience. A system that classifies independently of the operator removes this source of variance.

\textbf{Environmental consistency:} Dust, variable lighting, and shift-length fatigue progressively degrade classification reliability. Late-shift assessments showed more conservative borderline calls and slower sorting speeds than early-shift assessments. An automated system does not accumulate fatigue.

\textbf{Adulteration detection:} Quality managers described a specific practice: blending lower-grade grain into premium-variety lots, concentrated toward the inner \textit{boris} furthest from the periphery sample draw. The two-stage sampling protocol is specifically designed to catch this. An automated system must match the sensitivity of the most experienced assessors to avoid enabling rather than preventing quality fraud.
